\section{Claim Boundary}
\label{app:claims}

\begin{table}[h]
\centering
\caption{Current manuscript claim ledger.}
\small
\setlength{\tabcolsep}{3pt}
\begin{tabular}{p{0.16\linewidth}p{0.72\linewidth}}
\toprule
Status & Statement \\
\midrule
Supported & The 48-pair rejected-witness effect is small and heterogeneous under the frozen DeepSeek V2 design. \\
Supported & On the selected development stream, the historical strong First-Fail persistent state remains directionally useful across two frozen-state remeasurements. \\
Supported & Two fresh free-form updater realizations from reconstructed byte-identical First-Fail evidence do not reproduce the historical First-Fail advantage. \\
Supported, local & These observations localize persistent-state generation as a candidate bottleneck that cannot be reduced to evidence availability alone in the selected development case. \\
Not supported & Rejected failures are universally superior learning evidence. \\
Not supported & Progress proximity or diagnosticity defines a validated witness-selection law. \\
Not supported & Updater variance has been quantitatively separated from actor variance in the population. \\
Not supported yet & Typed compilation improves downstream utility or reproducibility. \\
Not supported yet & The mechanism generalizes to untouched E3, a second backbone, or a public benchmark. \\
\bottomrule
\end{tabular}
\end{table}

\section{Frozen Evidence Summary}
\label{app:evidence}

\paragraph{DeepSeek V2 Repair2.}
The frozen design contains 48 pairs, 96 learned states, and 1,728 held-out units. Mean MRW$-$WIN-C is $0.023148$, median $0.013889$, standard deviation $0.077970$, one-sided sign-flip $p=0.171875$, and bootstrap interval $[-0.018519,0.065972]$. The adjudicated verdict is \texttt{HOLD\_MRW\_UNDERPOWERED\_OR\_HETEROGENEOUS}.

\paragraph{Single-case witness selector.}
On \texttt{e1-tsr-00}, WIN-C/First-Fail/Progress-Fail/Progress-Contrast are 13/18, 17/18, 14/18, and 14/18. The frozen verdict is \texttt{S1\_SIGNAL\_FAIL\_STOP\_NO\_S2}.

\paragraph{Frozen-state stability.}
Two complete remeasurements give First-Fail versus WIN-C of 15/18 versus 14/18 and 16/18 versus 12/18. Verdict: \texttt{FIRST\_FAIL\_FROZEN\_STATE\_STABILITY\_PASS}.

\paragraph{Exact-evidence updater replay.}
All reconstructed learner-visible packet hashes match the historical First-Fail evidence. Fresh updater realizations give 15/18 versus 15/18 and 11/18 versus 12/18 against fresh WIN-C comparators. The frozen adjudication is that exact-evidence updater replication failed to reproduce the strong state, leaving state-generation realization as the candidate bottleneck.

\section{Prospective Stage Status}
\label{app:prospective}

M2, the deterministic G0--G3 state-semantic intervention, is frozen under an exactly-once Recovery V3 object. At this manuscript checkpoint the scientific result is not opened as a complete 72-unit analysis, and no partial effect is used in the paper. M3 is conditional on M2 passing its original gate and provider-era veto. M4, the fresh Evidence$\times$Generator bridge, is zero-provider proposal only. M5, untouched E3 confirmation, has no authority.

The paper therefore uses future-stage descriptions only as falsification contracts. They are not evidence and are not included in the abstract as positive method results.

\section{Bridge Suite Construction}
\label{app:suite}

The fresh bridge suite uses new controlled blocks 7--9 and has zero task-ID overlap with the earlier 378-task controlled suite. Blocks 7 and 8 provide update candidates; block 9 provides held-out candidates. From 162 deterministic candidates, hash-based selection freezes 96 update tasks in 12 eight-task streams, 12 SCREEN held-out tasks, 12 disjoint VALIDATION held-out tasks, 12 update integrity reserves, and 30 held-out integrity reserves. Each SCREEN/VALIDATION stream set contains one stream per failure family; each held-out panel contains two tasks per family. Reserved tasks may replace only a pre-execution file-integrity failure, never a model failure or unfavorable scientific outcome.

\section{Why the Paper Does Not Reuse the Old Temporal-Skill Manuscript}
The earlier E2 temporal-skill audit studied a different scientific object: comparator and integration-surface attribution for time-conditioned retrieval procedures. The present E2-R17 paper concerns persistent skill self-evolution from trajectory feedback. Reusing the old manuscript's abstract, results, or framing would conflate two independent research lines. This draft is therefore rebuilt from the current E2-R17 evidence ledger rather than incrementally editing the historical temporal-skill text.
